\section{Results}\label{sec:results}

The results organize around three findings. First, the simulated set-aside price effect is mostly within-auction: the within-auction share is \bneShareIntensNp~percent in non-pharma and \bneShareIntensPh~percent in pharma, and stays in the \withinShareLow--\withinShareHigh~percent range across cost-distribution estimators (Section~\ref{sec:bridge}). Second, endogenous SME entry is a partial offset, not the source of the markup: without it the realized price and welfare cost would be \latentShockRatioNp~percent larger in non-pharmaceuticals and \latentShockRatioPh~percent larger in pharmaceuticals. Third, the policy comparison between bidder exclusion and a \optPrefThreshold{} price preference is conditional on goods characteristics: in thick standardized non-pharma markets the preference rule welfare-dominates robustly; in thin heterogeneous pharma the ranking depends on how the post-policy entrant mix is modeled.

\subsection{Entry response}

Table~\ref{tab:v3_entry_rates} reports the panel. Pre-policy Preg\~ao has \bneNsmePreNp{} SMEs and \bneNnsPreNp{} non-SMEs per auction in non-pharma, \bneNsmePrePh{} and \bneNnsPrePh{} in pharma. Post-policy: \bneNsmePostNp{} SMEs and \nonSmePostNp{} non-SMEs in non-pharma, \bneNsmePostPh{} and \nonSmePostPh{} in pharma. SMEs roughly double; non-SMEs fall sharply. The set-aside therefore changes both the admissible pool and the equilibrium participation response.

\subsection{Structural decomposition of the set-aside}

Table~\ref{tab:v3_bne_decomp} gives the point estimates under the all-bidders regime and endogenous entry (the main specification). In non-pharma, $p_{S_1} = \bneMeanSoneNp$, $p_{S_2} = \bneMeanStwoNp$, $p_{S_3} = \bneMeanSthreeNp$; the within-auction shift is $+\bneEffectIntensNp$, the participation offset $\bneEffectEntryNp$, the net effect $+\bneEffectTotalNp$. Pharma follows the same decomposition with larger magnitudes: $p_{S_1} = \bneMeanSonePh$, $p_{S_2} = \bneMeanStwoPh$, $p_{S_3} = \bneMeanSthreePh$, total $+\bneEffectTotalPh$, intensive $+\bneEffectIntensPh$, participation $\bneEffectEntryPh$. The larger magnitudes track the thinner baseline SME pool (\bneNsmePrePh{} vs \bneNsmePreNp{}) and the heavier auction-level heterogeneity (ICC \iccDecompContrastPh{} vs \iccDecompContrastNp{} in Preg\~ao non-SME cells).

The within-auction share, $|p_{S_2} - p_{S_1}| / (|p_{S_2} - p_{S_1}| + |p_{S_3} - p_{S_2}|)$, is \bneShareIntensNp~percent in non-pharmaceuticals and \bneShareIntensPh~percent in pharmaceuticals in the main specification, \turnbullShareNp/\turnbullSharePh{} under the Turnbull NPMLE, and \strictShareNp/\strictSharePh{} under strict primitive invariance (Tables~\ref{tab:v3_sensitivity_fc} and~\ref{tab:v3_strict_invariance}). Across classes and specifications, the share sits in the \withinShareLow--\withinShareHigh~percent range. The set-aside mainly works by changing the order statistic inside the restricted auction; entry matters, but as a partial offset. Figure~\ref{fig:v3_decomposition} makes this visible.

\begin{figure}[!htb]
\centering
\includegraphics[width=0.98\textwidth]{../output/figures/fig_v3_decomposition.pdf}
\caption[Bayes-Nash counterfactual price decomposition.]{\textit{Bayes-Nash counterfactual price decomposition by class, under observed equilibrium entry.} Bars report the simulated second-order statistic $\bar{c}_{(2)}$ normalized by the item reference price, under three counterfactual scenarios. $S_1$ is the open regime at the Pre-period pool (baseline); $S_2$ imposes SME-only at the Pre-period SME pool (within-auction component isolated); $S_3$ is the SME-only counterfactual at the observed Post-period SME pool (full policy effect with endogenous entry). Annotations report the within-auction component ($S_2-S_1$), the participation-margin component ($S_3-S_2$), and the simulated total effect ($S_3-S_1$). Numbers from Table~\ref{tab:v3_bne_decomp}.}
\label{fig:v3_decomposition}
\end{figure}

\subsection{Endogenous entry as a partial offset}

Table~\ref{tab:v3_apv} makes the role of entry visible. V2 imposes the full set-aside while holding the SME pool at its pre-policy size; relative to V0, V2 is \latentShockTimesNp~percent of V0 in non-pharmaceuticals and \latentShockTimesPh~percent in pharmaceuticals---the fixed-pool simulated effect exceeds the realized one by \latentShockRatioNp/\latentShockRatioPh~percent. Both terms in the welfare formula of Section~\ref{sec:welfare} load on this margin, so a fixed-pool calculation would materially overstate the policy burden. Under V3 (10 percent price preference) the analogous fixed-pool-vs-realized comparison delivers an attenuation close to zero since non-SMEs remain in the auction.

\subsection{Bridges, CIs, and entry costs}\label{sec:bridge}

The 2--3$\times$ gap between the structural counterfactual ($+\bneEffectTotalNp/+\bneEffectTotalPh$) and the DiD-implied ratio (\didImpliedRatio~percent on absolute prices) is decomposed into four mechanical sources (sample restriction, UH cleaning, functional form, conditioning set) summing to $\bridgeSumNp/\bridgeSumPh$ on the $p^{\mathrm{ref}}$ normalization, with residuals $\bridgeResidNp/\bridgeResidPh$ from non-linear interactions (Table~\ref{tab:did_structural_bridge}; full grid in Table~\ref{tab:v3_decomp_grid}). Functional form is the dominant source in non-pharma ($\bridgeFuncNp$); UH cleaning in pharma ($\bridgeUhPh$). The gap is predictable in direction and magnitude under correct specification, not a misspecification flag. A cluster bootstrap at the auction level ($\bootReps$ replicates, Table~\ref{tab:v3_bootstrap_ci}) gives 95 percent CIs on $p_{S_3} - p_{S_1}$ of $[\bootCILoNp, \bootCIHiNp]$ non-pharma and $[\bootCILoPh, \bootCIHiPh]$ pharma; both exclude zero. Within-auction-share CIs are $[\bootShareCILoNp\%, \bootShareCIHiNp\%]$/$[\bootShareCILoPh\%, \bootShareCIHiPh\%]$, with the intensive component larger throughout. Zero-profit implied entry costs (Table~\ref{tab:v3_entry_cost}, \entryCostMin--\entryCostMax{} per bidder, non-SME-to-SME ratio of \entryCostRatioNp-to-1) provide a calibration check on the Poisson primitives rather than a free-standing structural finding; the Online Appendix expands on each diagnostic.

\subsection{Policy comparison and goods characteristics}\label{sec:pareto}

The policy comparison follows directly from the structural decomposition: V3 leaves non-SMEs in the auction, V0 removes them entirely, and the within-auction component prices the difference. Under a \optPrefThreshold{} SME price preference, SME bids are scored with a \optPrefThreshold{} discount for winner selection but the government pays the actual bid; in non-pharma the price shift is $\optPrefShiftNp$ relative to the open regime, in pharma $\vThreeShiftPh$. The preference-grid simulation ($\prefGridB$ draws, Table~\ref{tab:v3_preference_grid}) gives $\vThreePrefGrid$ in both classes within Monte Carlo noise of the APV estimates. The price effect is essentially zero because non-SMEs continue to discipline the second-order statistic even when an SME wins.

What differs across classes is the stability of the ranking against the full set-aside. In non-pharma, the preference rule dominates under both the main equilibrium-selection specification and the strict-invariance benchmark. In pharma, the preference rule dominates under the main specification but the ranking reverses under strict invariance, where the full set-aside's total price effect rises to $+\siDeltaTotalPh$ and the implied welfare weight falls to \welfWeightSiPh. The bifurcation is diagnostic of market structure: in thicker, more standardized markets, bidder exclusion is robustly dominated by a moderate preference; in thinner, more heterogeneous markets, the ranking depends on the identity of the induced entrants---a property the framework identifies but a single-platform exercise cannot resolve.
